\documentclass[a4paper, 12pt]{article}

\usepackage{utils}

\renewcommand*{\today}{04 March 2026}

\begin{document}

\hotbox{Analyse 4}{CM 9}{\today}

\begin{proposition}{Formule de changement de variable}{}
    On prend la série entière $f(z) := \sum\limits_{n \geq 0} a_n z^n$ de rayon de convergence $R > 0$ et $|z| < R$.

    Alors pour tout $\alpha \in \R$, pour tout $m \in \Z^*$, on a
    $$
    \sum_{n \geq 0} a_n \alpha^n z^{nm} = f(\alpha z^m)
    $$
    dés que $|\alpha z^m| < R \iff z \leq \sqrt[m]{\dfrac{R}{|\alpha|}}$
\end{proposition}

\begin{proposition}{}{}
    Pour tout $z \in \mathbb{C}$, $e^z := \sum_{n \geq 0} \dfrac{z^n}{n!}$.

    Alors
    $$
    e^{a + b} = e^a e^b
    $$
    pour tout $a, b \in \mathbb{C}$.
\end{proposition}

\begin{demonstration}
    On a
    $$
    e^a = \sum_{n \geq 0} \dfrac{a^n}{n!} = \sum_{n \geq 0} 1^n \dfrac{a^n}{n!} = (z \mapsto e^{a z})(1)
    $$
    pareil pour $e^b$.
    $$
    e^b = (z \mapsto e^{b z})(1)
    $$

    On fait le produit des deux séries
    $$
    e^a e^b = \left(\sum_{n \geq 0} 1^n \dfrac{a^n}{n!} \right)\left( \sum_{n \geq 0} 1^n \dfrac{b^n}{n!} \right) = \sum_{n \geq 0} c_n 1^n
    $$
    avec
    \begin{align*}
        c_n := \sum_{i + j = n} \dfrac{a^i}{i!} \dfrac{b^j}{j!} &= \dfrac{b^n}{n!} + \dfrac{ab^{n-1}}{(n-1)!1!} + \cdots + \dfrac{a^n}{n!} \\
        &= \sum_{k = 0}^n \dfrac{b^{n-k}a^k}{(n-k)!k!} = \dfrac{1}{n!} \sum_{k = 0}^n \binom{n}{k} a^k b^{n-k} = \dfrac{(a + b)^n}{n!}
    \end{align*}

    et on a aussi

    $$
    e^a \times e^b = \sum_{k \geq 0} \dfrac{(a + b)^k}{k!} 1^k = e^{a + b}
    $$
\end{demonstration}

\subsection{Développabilité en série entière}

On ne considère que $z \in \R$.

\begin{definition}
    Soit $z_0 \in \R$, $\bar{\R^+} = \R^+ \cup \{+\infty\}$ et $R \in \bar{\R^+}$.

    $g$ un fonction définie sur $\R$ sera \textbf{dévéloppable} en série entière sur $]z_0 - R, z_0 + R[$ si
    $$
    \forall z \in ]z_0 - R, z_0 + R[, \quad g(z) = \sum_{n \geq 0} a_n (z - z_0)^n
    $$
    cette somme s'appellera la \textbf{développement en série entière} de $g$ au voisinage de $z_0$.
\end{definition}

\begin{proposition}{}{}
    La fonction $g$ est développable en série entière sur $]z_0 - R, z_0 + R[$ si et seulement si
    la fonction $z \mapsto g(z + z_0)$ est développable en série entière sur $]-R, R[$ en $z = 0$.
\end{proposition}

\begin{theoreme}{}{}
    Soit $f$ une fonction d.s.e autour de 0 et $R > 0$ son rayon de convergence telle que $\forall z \in ]-R, R[, f(z) = \sum\limits_{n \geq 0} a_n z^n$.

    Alors $f$ est indéfiniment dérivable sur $]-R, R[$ et on dérive terme à terme $f'(z) = \sum\limits_{n \geq 1} n a_n z^{n-1}$.

    Plus généralement
    $$
    f^{(l)}(z) = \sum_{n \geq l} \prod_{k = 0}^{l-1} (n - k) a_n z^{n-l}
    $$
    et de plus
    $$
    \int_0^z f(t) dt = \sum_{n \geq 0} \dfrac{a_n}{n + 1} z^{n + 1}
    $$
\end{theoreme}

\begin{demonstration}
    On a
    $$
    \lim_{n \to +\infty} \sqrt[n]{n} = \lim_{n \to +\infty} n^{\frac{1}{n}} = \lim_{n \to +\infty} e^{\frac{\ln(n)}{n}} = 1
    $$
    et donc $\lim_{n \to +\infty} \sup \sqrt[n]{n |a_n|} = \lim_{n \to +\infty} \sup \sqrt[n]{|a_n|}$ donc les deux séries
    $\sum_{n \geq 0} na_n z^n$ et $\sum_{n \geq 0} a_n z^n$ ont le même rayon de convergence $R$.

    D'apres la proposition ?? on a la convergence uniforme de $\sum_{n \geq 0} a_n z^n$ sur $[-r, r]$ pour tout $r \in ]0, R[$.
    De même pour $\sum_{n \geq 0} na_n z^n$.

    On peut appliquer le théorème de dérivation d'une série de fonctions pour conclure que $f$ est dérivable et que $f'(z) = \sum\limits_{n \geq 1} n a_n z^{n-1}$.

    \vspace{0.5cm}
    Par récurrence, on trouve la formule pour $f^{(l)}$.
    \vspace{0.5cm}

    Pour l'intégrale, on applique le théorème de dérivation à $\int_0^z f(t) dt$.
\end{demonstration}

\begin{corollaire}{}{}
    Si $f$ est d.s.e sur $]-R, R[$ avec $R > 0$.

    Alors
    $$
    f(z) = \sum_{n \geq 0} \dfrac{f^{(n)}(0)}{n!} z^n
    $$
\end{corollaire}

\begin{remarque}
    On reconnait un développement de Taylor de $f$ au voisinage de 0.
\end{remarque}

\begin{demonstration}
    On fait une récurrence sur $n$
\end{demonstration}

\begin{definition}
    Si $f$ est indéfiniment dérivable sur $]-R, R[$, alors le developpement,
    $$
    f(z) = \sum_{n \geq 0} \dfrac{f^{(n)}(0)}{n!} z^n
    $$
    s'appelle la série de Taylor de $f$ au voisinage de 0.
\end{definition}


\end{document}